%% La contrainte normale sigma = F / S, en trois temps :
%% 1. la barre tendue ; 2. la coupure imaginaire et l'isolement d'un tronçon ;
%% 3. la section coupée, où l'effort normal N se répartit uniformément.
\documentclass[tikz,border=4pt]{standalone}
\usepackage{liaisons}
\begin{document}
\begin{tikzpicture}[x=1cm,y=1cm,
  force/.style={-{Stealth[length=6pt]}, line width=1.2pt, solideA},
  petite/.style={-{Stealth[length=3.6pt]}, line width=0.7pt, solideA},
  etape/.style={-{Stealth[length=5pt]}, line width=0.7pt, black!55},
  titre/.style={font=\small, anchor=north, align=center}]

%% ==================== 1. la barre tendue ============================
\begin{scope}[shift={(0,0)}]
  \draw[line width=1pt, solideB, fill=solideB!10, rounded corners=2pt] (-0.8,-0.3) rectangle (0.8,0.3);
  \draw[force] (0.85,0) -- (1.45,0) node[font=\small, text=solideA, anchor=south, inner sep=2pt] {$F$};
  \draw[force] (-0.85,0) -- (-1.45,0) node[font=\small, text=solideA, anchor=south, inner sep=2pt] {$F$};
  \node[titre] at (0,-0.75) {\textbf{1.} la pièce\\ est tendue};
\end{scope}
\draw[etape] (1.70,0) -- (2.15,0);

%% =============== 2. la coupure imaginaire ===========================
\begin{scope}[shift={(3.75,0)}]
  \draw[line width=1pt, solideB, fill=solideB!10, rounded corners=2pt] (-0.8,-0.3) rectangle (0,0.3);
  \draw[line width=0.7pt, draw=black!45, fill=black!10, rounded corners=2pt] (0,-0.3) rectangle (0.8,0.3);
  \draw[line width=1.1pt, solideD, dash pattern=on 3pt off 2pt] (0,-0.5) -- (0,0.5);
  \draw[force] (-0.85,0) -- (-1.45,0) node[font=\small, text=solideA, anchor=south, inner sep=2pt] {$F$};
  \node[font=\small, anchor=south, inner sep=2pt, text=black!60] at (0.42,0.34) {supprimé};
  \node[titre] at (-0.1,-0.75) {\textbf{2.} coupure imaginaire :\\ on isole un tronçon};
\end{scope}
\draw[etape] (5.15,0) -- (5.60,0);

%% =============== 3. la section coupée ===============================
\begin{scope}[shift={(6.6,0)}]
  \hachures{(0,0) ellipse (0.3 and 0.62)}
  \foreach \y in {-0.46,-0.23,0,0.23,0.46}
    { \pgfmathsetmacro\dx{0.3*sqrt(1-(\y/0.62)^2)}
      \draw[petite] (\dx,\y) -- ++(0.42,0); }
  \fill (0,0) circle (1.6pt);
  \node[font=\small, anchor=east, fill=white, inner sep=1pt] at (-0.06,0.02) {$G$};
  \node[font=\small, text=solideA, anchor=west, inner sep=2pt] at (0.76,0) {$N$};
  \node[font=\small, anchor=south, align=center] at (0.3,0.68) {répartition\\ uniforme};
  \node[titre] at (0.2,-0.75) {\textbf{3.} l'effort $N$ se\\ répartit sur $S$};
\end{scope}

%% ------------------------- cartouche --------------------------------
\node[draw=black!45, fill=black!3, rounded corners=3pt, inner sep=5pt, anchor=north,
      font=\small, align=center] at (3.1,-1.72)
     {$\sigma = \dfrac{F}{S}$ \quad $F$ en N, $S$ en mm$^2$, $\sigma$ en N/mm$^2$ = MPa\\[2pt]
      section circulaire de diamètre $d$ : $S = \dfrac{\pi d^2}{4}$};
\end{tikzpicture}
\end{document}
